% =====================================================================
% 2. Referencial histórico
% =====================================================================
\section{Referencial histórico}
\label{sec:historia}

\subsection{Tabelas, logaritmos e interpolação}

Durante séculos, calcular significou consultar e completar tabelas. Efemérides astronômicas, tábuas de tiro, tabelas de juros e funções transcendentes eram produzidas à mão, com conferência dupla, porque um dígito errado se propagava para todos os usuários da tábua. A interpolação era a ferramenta que extraía um valor intermediário sem refazer a conta inteira. As fórmulas de diferenças finitas, mais tarde associadas a Newton e a Lagrange, nasceram dessa oficina, não de uma preocupação abstrata com polinômios \citep{goldstine1972,burden2016}.

A invenção dos logaritmos, no início do século XVII, mudou o custo da multiplicação. Uma multiplicação virava soma de logaritmos, e uma raiz virava divisão do logaritmo. As tábuas de logaritmos passaram a ser infraestrutura de engenheiros, navegadores e astrônomos. A régua de cálculo é a versão analógica dessa ideia: escalas logarítmicas deslizantes fazem a soma dos logaritmos por um deslocamento mecânico. William Oughtred formulou a régua no século XVII, e o instrumento permaneceu na bancada de engenheiros até as calculadoras científicas eletrônicas da década de 1970 \citep{hopp1999,williams1997}. A precisão era de poucos algarismos, suficiente para uma estimativa de projeto e insuficiente para uma tabela fina. Essa limitação já anunciava um tema deste trabalho: o instrumento determina o erro tanto quanto a fórmula.

Havia também tabelas de números aleatórios. Antes de geradores computacionais confiáveis, simulação, estatística e criptografia usavam dígitos impressos. O caso mais conhecido é \emph{A Million Random Digits with 100,000 Normal Deviates}, publicado pela RAND Corporation em 1955, com um milhão de dígitos obtidos por um dispositivo eletrônico semelhante a uma roleta \citep{rand1955,lecuyer2017}. A curiosidade mostra o quanto o cálculo numérico dependia de objetos físicos, e não apenas de teoremas.

\subsection{De Newton a Runge e Kutta}

Isaac Newton deixou, no século XVII, o método das diferenças finitas e um procedimento iterativo para raízes de polinômios. Joseph Raphson publicou em 1690 uma forma mais direta desse procedimento, que a tradição posterior consolidou como método de Newton-Raphson, já com a derivada explícita. O uso imediato era corrigir tábuas e resolver equações que não se deixavam inverter em forma fechada \citep{burden2016,franco2006}.

Leonhard Euler publicou em 1768, no primeiro volume de \emph{Institutionum calculi integralis}, o que hoje chamamos de método de Euler para equações diferenciais: avançar a solução com a derivada do ponto atual multiplicada pelo passo \citep{euler1768,hairer1993}. A ideia é a definição de derivada lida ao contrário. Euler também desenvolveu a fórmula de Euler-Maclaurin, que liga somas a integrais e está na origem de várias correções de quadratura.

Joseph-Louis Lagrange reorganizou a interpolação polinomial e o cálculo de diferenças numa linguagem que ainda reconhecemos. Carl Friedrich Gauss, no começo do século XIX, deu forma sistemática à eliminação para os sistemas lineares dos mínimos quadrados da astronomia, e construiu a quadratura que leva seu nome. Um procedimento equivalente de eliminação já aparecia nos \emph{Nove capítulos da arte matemática}, na China antiga; o que Gauss fez foi torná-lo ferramenta regular da ciência aplicada \citep{goldstine1972,franco2006}. No século XIX, Carl Gustav Jacobi propôs um método iterativo para os sistemas lineares que surgem nos mínimos quadrados. A iteração de Jacobi, ao lado da de Gauss-Seidel, é o antepassado dos métodos usados hoje em sistemas esparsos de grande porte \citep{burden2016,ruggiero1996}.

A linha que este trabalho segue de perto começa no fim do século XIX. Carl Runge publicou em 1895 um método de integração que avaliava a função em estágios internos, em vez de exigir derivadas totais de ordem alta \citep{runge1895}. Karl Heun, em 1900, sistematizou a comparação com a série de Taylor e construiu métodos de ordem mais alta, entre eles o corretor que a disciplina chama de Euler melhorado \citep{heun1900}. Wilhelm Kutta, em 1901, deu a forma geral dos métodos explícitos de estágios e apresentou o método clássico de quarta ordem, com quatro avaliações da função por passo \citep{kutta1901}. Butcher e Wanner \citep{butcher1996} reconstruíram essa sequência e mostraram como as condições de ordem, a estabilidade linear e os métodos implícitos saíram dali no século seguinte. O algoritmo que um programa executa em poucas linhas foi, na origem, um substituto engenhoso da série de Taylor para quem calculava à mão.

\subsection{Instrumentos de cálculo}

O ábaco, em versões mesopotâmicas, romanas e chinesas, é o registro antigo de um princípio que a computação não abandonou: separar o estado (as contas) da regra (o movimento). A pascalina de Blaise Pascal, concebida em 1642, mecanizou a soma e o transporte. A calculadora de Gottfried Wilhelm Leibniz, o \emph{Stepped Reckoner}, estendeu a máquina às quatro operações por meio do cilindro dentado de Leibniz. No mesmo ambiente intelectual, Leibniz descreveu a aritmética binária em 1703 \citep{leibniz1703}. O texto é filosófico e matemático, não um projeto de computador, mas a representação que os circuitos digitais adotaram está formulada ali \citep{goldstine1972,williams1997}.

Charles Babbage deu o passo seguinte ao separar a produção de tabelas de um operador humano. A Máquina de Diferenças, projetada a partir de 1822, avaliava polinômios pelo método das diferenças finitas: depois de algumas colunas iniciais, só resta somar. A Máquina de Diferenças nº~2, desenhada entre 1847 e 1849, só foi construída pelo Science Museum de Londres na década de 1990, sob a coordenação de Doron Swade, e funcionou. O desenho de Babbage era viável; faltaram, no século XIX, dinheiro, tolerância mecânica e tempo político \citep{swade2001}.

A Máquina Analítica, descrita por Babbage a partir de 1837, já não era uma máquina de uma fórmula só. Tinha moinho (\emph{mill}), depósito (\emph{store}), cartões de operação e cartões de variável, inspirados no tear de Jacquard, e previa repetição e desvio. Não foi concluída. Ainda assim, é o primeiro projeto reconhecível de um computador de programa externo \citep{swade2001,goldstine1972}.

\subsection{Lovelace, Babbage e a ideia de programa}

Em 1843, Augusta Ada King, Condessa de Lovelace, traduziu para o inglês o artigo de Luigi Menabrea sobre a Máquina Analítica e acrescentou notas mais longas do que o próprio artigo \citep{menabrea1843}. A Nota G descreve, passo a passo, o cálculo dos números de Bernoulli na máquina. É o documento que a história da computação trata como o primeiro programa publicado. Lovelace também percebeu que a máquina manipulava símbolos segundo regras, e que o conteúdo desses símbolos não precisava ser apenas número no sentido aritmético estreito. A distinção entre o mecanismo e a sequência de operações é a distinção entre hardware e programa.

Babbage queria automatizar o cálculo que as tabelas exigiam. Lovelace descreveu a automação como uma linguagem de procedimentos. Os dois lados continuam no trabalho de quem implementa um método numérico: há a fórmula, e há a ordem das avaliações, os testes de parada e a representação dos dados. Sem essa segunda camada, o método de Kutta permanece uma receita de papel.

\subsection{Computadores eletrônicos e a arquitetura de von Neumann}

Entre a engrenagem e a válvula houve uma geração eletromecânica. O Z3 de Konrad Zuse, em 1941, operava com relés, aritmética binária de ponto flutuante e programa em fita perfurada. É amplamente reconhecido como o primeiro computador programável automático em funcionamento \citep{rojas1997}. O Harvard Mark I, de Howard Aiken, seguia a mesma família eletromecânica, em grande escala. Essas máquinas são a transição, não um detalhe: o cálculo deixou de ser um gesto da mão antes de deixar de ser mecânico.

O ENIAC, construído na Universidade da Pensilvânia por J. Presper Eckert e John Mauchly, com participação de Herman Goldstine, foi apresentado em 1946 para calcular tabelas balísticas. Era eletrônico, decimal, e a programação original era feita por cabos e chaves. Programar era reconfigurar o hardware. Em 1948 o ENIAC recebeu um modo rudimentar de programa armazenado \citep{goldstine1972}. Os cartões perfurados, herdeiros do sistema de Herman Hollerith para o censo americano de 1890, continuaram por décadas como suporte de programa e de dados. A separação entre escrever um algoritmo e executá-lo, hoje trivial, foi uma conquista material.

O texto que fixou a organização depois dominante é o \emph{First Draft of a Report on the EDVAC}, de John von Neumann, de 1945 \citep{vonneumann1945}. Unidade aritmética, memória que guarda instruções e dados, unidade de controle e órgãos de entrada e saída: a chamada arquitetura de von Neumann. O próprio Goldstine \citep{goldstine1972} registra que a ideia de programa armazenado foi elaborada no grupo da Moore School, com Eckert e Mauchly, e não por uma pessoa só. O primeiro programa armazenado executado numa máquina eletrônica correu no Manchester Baby, de Frederic Williams e Tom Kilburn, em 21 de junho de 1948. O EDSAC, de Maurice Wilkes, em Cambridge, entrou em serviço regular em 1949 e tornou o modelo utilizável para cálculo científico \citep{williams1997}. A partir daí, um método numérico podia ser texto, copiado e reexecutado, e não mais um painel de cabos.

\subsection{Aritmética binária, ponto flutuante e o limite do hardware}

O Z1 e o Z3 de Zuse já usavam ponto flutuante binário na passagem da década de 1930 para a de 1940 \citep{rojas1997}. A padronização que os computadores atuais seguem veio bem depois: o padrão IEEE 754, publicado em 1985 e revisado em 2008 e 2019, fixou formatos, arredondamentos e exceções \citep{ieee2019}. Um número em precisão dupla (\emph{binary64}) tem um bit de sinal, 11 de expoente e 52 de fração. A precisão simples (\emph{binary32}) tem 8 bits de expoente e 23 de fração. O épsilon de máquina, distância relativa entre 1 e o próximo representável, vale cerca de \(2{,}22\times 10^{-16}\) na precisão dupla e \(1{,}19\times 10^{-7}\) na precisão simples.

Essa representação não é um detalhe de implementação. Ela introduz erro de arredondamento em quase toda operação, faz somas deixarem de ser exatamente associativas e cria cancelamento quando se subtraem quantidades próximas \citep{goldberg1991,overton2001,higham2002}. Overflow, underflow e perda de dígitos significativos são limites físicos do formato, não falhas do método matemático. Numa integração longa, o erro de truncamento do método cai com o passo, enquanto o erro de arredondamento tende a crescer, porque há mais operações. A Seção~\ref{sec:precisao} mede esse cruzamento no circuito RLC.

O hardware também limita a estabilidade por outro caminho. Um método explícito exige um passo pequeno quando a dinâmica é rápida. Quanto maior a frequência do circuito, menor o passo em tempo absoluto para percorrer o mesmo número de períodos. O algoritmo pode ser o mesmo; o tempo de máquina não é.

\subsection{Linguagens e bibliotecas numéricas}

Enquanto o programa era escrito em código de máquina, o método numérico e o endereço dos registradores se misturavam. Fortran, desenvolvido na IBM por uma equipe liderada por John Backus, separou as duas coisas. O primeiro compilador foi entregue em 1957 para o IBM 704. A linguagem foi desenhada para traduzir fórmulas para código eficiente, e tornou-se o veículo do cálculo científico durante décadas \citep{backus1979}.

As bibliotecas seguiram essa vocação. Os Basic Linear Algebra Subprograms (BLAS), de 1979, padronizaram as operações de vetor e, nas extensões posteriores, de matriz \citep{lawson1979}. O LAPACK, a partir de 1992, construiu sobre o BLAS os métodos de álgebra linear densa usados por quase todo software científico \citep{anderson1999}. A FFTW mostrou que uma transformada tão clássica quanto a de Fourier ainda ganhava desempenho com planejamento automático e atenção à hierarquia de memória \citep{frigo2005}. O cálculo numérico de alto desempenho, hoje, é tanto matemática quanto organização de dados.

MATLAB nasceu no fim da década de 1970, quando Cleve Moler quis uma interface interativa para as bibliotecas LINPACK e EISPACK no ensino. A MathWorks foi fundada em 1984, e o programa tornou-se o ambiente padrão de muita engenharia \citep{moler2004}. O Scilab teve origem no ambiente Basile, do INRIA, e a versão 1.1 foi distribuída em 1994 como software científico aberto \citep{scilab2013,gomez1999}. No ecossistema Python, o SciPy foi lançado em 2001 e o NumPy 1.0 em 2006, reunindo o trabalho anterior do Numeric e do Numarray \citep{virtanen2020,harris2020}. Julia, apresentada publicamente em 2012, foi desenhada para que o código numérico de alto nível compilasse para desempenho próximo do de linguagens de sistemas \citep{bezanson2017}. A recomendação da disciplina, que admite Scilab, Python, MATLAB, Octave, C/C++, Java, Julia ou Pascal, reflete exatamente essa diversidade. A escolha deste trabalho é Python, com os métodos escritos no próprio código.

\subsection{O cálculo numérico na engenharia contemporânea}

O mesmo núcleo reaparece em áreas que parecem distantes da tábua de logaritmos. O treinamento de redes neurais é otimização numérica, em geral uma variação de gradiente, sobre álgebra linear de ponto flutuante. Formatos de 16 bits e o \emph{bfloat16} existem para reduzir memória e tempo; o preço é o retorno dos cuidados clássicos com escala, cancelamento e condicionamento \citep{higham2002,ieee2019}. Simulação científica de estruturas, fluidos e circuitos continua a integrar equações diferenciais e a resolver sistemas lineares esparsos. O processamento de sinais e as comunicações digitais dependem de amostragem, filtragem e transformadas rápidas; um estágio ressonante analógico é o protótipo contínuo de uma seção de filtro digital \citep{oppenheim1997,lathi2009}.

Na criptografia, o miolo dos algoritmos modernos é discreto, mas a análise numérica entra na geração de aleatoriedade e em ferramentas como o algoritmo LLL de redução de reticulados, cuja implementação clássica usa ortogonalização em ponto flutuante \citep{lenstra1982,lecuyer2017}. Na computação paralela, a álgebra linear densa se distribui com naturalidade; o avanço no tempo de um método explícito de passo único, ao contrário, é uma cadeia de dependências. Paralelizar a avaliação de um sistema grande, com muitos estados, é viável. Paralelizar o próprio tempo exige métodos específicos, como os esquemas parareais, e não sai de graça de um laço de Runge-Kutta.

Os computadores analógicos fecham o arco. Construídos com amplificadores operacionais, representavam variáveis por tensões e realizavam soma, ganho e integração por circuitos. Um integrador analógico resolvia equações diferenciais em tempo contínuo, com aplicação forte em controle, aeronáutica e telecomunicações, antes de o digital prevalecer \citep{jung2005,chm}. A precisão era limitada por \emph{offset}, deriva térmica e tolerância dos componentes. O método digital troca esse erro físico pelo erro de truncamento e de arredondamento, que pode ser medido e reduzido com passo, ordem e formato. O experimento deste trabalho é essa troca, feita num circuito que um computador analógico também saberia integrar.
